\documentclass{report}
\usepackage{amsmath,amssymb}
\setlength{\parindent}{0mm}
\setlength{\parskip}{1em}
\begin{document}
\begin{center}
\rule{6in}{1pt} \
{\large Bram Reps \\
{\bf Communication avoiding strategies for linear solvers in a space-filling curve framework}}

Department of Mathematics and Computer Science \\ University of Antwerp \\ Middelheimlaan 1 \\ B-2020 Antwerp \\ Belgium \\ Intel ExaScience Lab \\ Kapeldreef 75 \\ B-3001 Leuven \\ Belgium
\\
{\tt bram.reps@ua.ac.be}\\
Bart Verleye\\
Tobias Weinzierl\\
Dirk Roose\\
Pieter Ghysels\\
Wim Vanroose\end{center}

Iterative methods such as Krylov subspace and multigrid methods are
typically used to solve large linear systems that arise from the
discretization of a partial differential equation (PDE). As these
equations mainly connect neighboring points in space, resulting in very
sparse systems, a data structure that is based on space-filling curves
can be exploited to develop numerical solvers, since the associated data
mappings carry specific locality properties \cite{B13}.

The cost of an iterative method is largely determined by the efficient
use of two basic calculations: sparse matrix-vector multiplications
(SpMV) for stencil operations, and vector multiplications (dot products)
for computing projections and norms. Both of these common steps in a
numerical algorithm require time-consuming data communication. The result
of a SpMV is a vector and communication is locally restricted to a few
neighboring points in space, as given by the sparsity pattern of the
matrix, whereas a dot product results in a scalar that depends on all
points in space.

We study a framework where a multidimensional numerical grid is traversed
in a particular order defined by a one-dimensional space-filling curve.
In each cell along this curve only local grid data, belonging to the cell
and its adjacent vertices, is manipulated as specified by the iterative
method, before moving on to the next cell in line. This approach is
particularly interesting for the sequentialization of adaptive grids,
since the curve can recursively be extended on-the-fly in refining
regions and the underlying data is typically stored in a spacetree. In
addition, the parallellization of a solver is straightforward if the
space-filling curve is simply cut and divided over different processors.

Although the data dependency for the SpMV is usually relatively local in
space, an exchange of information between cells is unavoidable in between
computations. This can easily be invoked by simply running extra grid
traversals dedicated to data communication through mutual vertices;
however, it dramatically increases the runtime of a programming code,
especially when communicating cells are handled by different parallel
processors. The result of a dot product is stored in a state variable,
accessible from all cells, and needs to be communicated back and forth
between parallel processors at the end of a traversal.

It fits better to the state-of-the-art in high performance computing if
the two phases of computation and communication are combined as much as
possible. This means that from the point of view of space-filling curves,
the current tendency of communication avoiding and hiding programming
\cite{H10,GAMV12} naturally translates into the basic concept of
minimizing the number of grid traversals by postponing and bundling the
communication of SpMV and dot product results.

To elucidate these strategies, we present a parallel implementation of
both the Conjugant Gradient (CG) method \cite{HS52} and an additive
multigrid method that respects the idea of minimal grid traversals and
illustrate them in the space-filling curve framework Peano
\cite{W09,W12}. The necessary communication of dot product and SpMV
results, is aggregated in few instances of the algorithm's main loop,
requiring e.g.\ only one grid traversal per CG-iteration.

\begin{thebibliography}{10}
\bibitem{B13}
M.~Bader.
\newblock {\em Space-Filling Curves}.
\newblock Springer-Verlag Berlin Heidelberg, 2013.

\bibitem{GAMV12}
P.~Ghysels, T.~J. Ashby, K.~Meerbergen, and W.~Vanroose.
\newblock Hiding global communication latency in the {GMRES} algorithm on
massively parallel machines.
\newblock {\em SIAM Journal on Scientific Computing}, 2012.

\bibitem{HS52}
M.~Hestenes and E.~Stiefel.
\newblock Methods of {C}onjugate {G}radients for solving linear systems.
\newblock {\em Journal of Research of the National Bureau of Standards}, 49(6),
December 1952.

\bibitem{H10}
M.~Hoemmen.
\newblock {\em Communication-avoiding Krylov subspace methods}.
\newblock PhD thesis, University of California, 2010.

\bibitem{W09}
T.~Weinzierl.
\newblock {\em A Framework for Parallel PDE Solvers on Multiscale Adaptive
Cartesian Grids}.
\newblock PhD thesis, Technische Universit�t M�nchen, 2009.

\bibitem{W12}
T.~Weinzierl et~al.
\newblock Peano -- a {F}ramework for {PDE} {S}olvers on {S}pacetree {G}rids,
2012.
\newblock http://www.peano-framework.org.
\end{thebibliography}


\end{document}
